
Weight space learning (WSL) methods that exploit permutation equivariance achieve high within-architecture performance on convolutional neural network (CNN) benchmarks but generalize poorly across architecture families. This paper investigates orbit-based positional encodings (orbit-PE) derived from the input/output channel permutation group as a candidate for architecture-agnostic weight tokenization. Three empirical measurements are reported. First, the input/output channel permutation group is confirmed as a functionally exact symmetry for Conv2d, Linear, and MultiheadAttention layers: |Delta acc| = 0.000000 across 4,500 permutation runs on CNN Zoo and Transformer Zoo checkpoints. Second, orbit-PE is computable for all layer types with a mean overhead of 1.167x relative to sequential positional encoding, using a unified codebase with no architecture-specific branches. Third, variance decomposition of 1,000 CNN Zoo models over 50 training epochs reveals a statistically robust layer-type stratification: the permutation-orbit variance fraction is 0.637 for Conv2d layers and 0.133 for Linear (fully-connected) layers, with an overall ratio of 0.3479 +/- 0.0536 — substantially below the pre-specified threshold of 0.60. The overall gate failure blocked planned cross-architecture training experiments. An ancillary finding is that the permutation variance fraction decreases monotonically during training, from approximately 0.49 at epoch 0 to approximately 0.28 at epoch 50. These results indicate that permutation orbit encoding is insufficient as the sole positional encoding for cross-architecture WSL and motivate a layer-type-specific hybrid encoding combining permutation orbit-PE for convolutional layers with GL-invariant trace features for linear and attention layers.

